\scanpage{112}
et ici, dans un sens très fort (qu'il vaudrait la peine peut-être d'expliciter) $P_C\times m_C$ s'identifie à une trace de voisinage tubulaire, \deleted{et $P_C$ serait le bord}

Les cas les plus intéressants sont ceux où $G'=G^{\unsure{\natural|\sharp}}$ (générisation maximale de $G$), ainsi que ceux où $G$ est de dim. modulaire $0$ ou $1$, \deleted{(\unsure{ou le} \unsure{fonct|front}} et où $G'$ est de dim modulaire $\leq2$ --- notamment les cas d'une générisation immédiate (correspondant au cas où $C$ est réduit à un élément).

Cas où la dim modulaire $d'(G)$ de $G$ est nulle.
\margin{ie les sommets de $G$ \unsure{sont} de genre $0$, d'ordre $3$.}
Alors $X$ est défini à isom. unique près par la donnée de $G$, grâce au fait qu'une droite projective munie d'un \deleted{$\unsure{B}$} $I$-marquage, avec $\operatorname{card}I=3$, est défini en termes de $I$ à isom. canonique près. Donc la multiplicité $M_G$ est la multiplicité ponctuelle. Considérons \deleted{le groupe} la droite projective $\mathbb P_{\mathbb C}$ avec les
\mspage{111}
